\chapter{Diagrammi dei casi d'uso}

Nel Capitolo 5 abbiamo visto come si \textbf{scoprono} i requisiti (interviste, personas, storie e scenari, mock-up) e i quattro
approcci per \textbf{specificarli}. Uno di questi approcci usa notazioni \textbf{semi-formali}: in questo capitolo vediamo la più
diffusa per i requisiti funzionali, i \textbf{casi d'uso} (\textit{use cases}) di UML, con il loro diagramma e la loro descrizione
testuale.

\begin{figure}[H]
\centering
\begin{tikzpicture}[font=\scriptsize]
  \node[card=ph1, text width=4.3cm, align=left, minimum height=2.3cm] (el) {\textbf{\small Raccolti tramite}\\[3pt]
    \textbullet\ interviste con gli stakeholder\\ \textbullet\ personas\\ \textbullet\ storie e scenari};
  \node[card=ph4, text width=4.3cm, align=left, minimum height=2.3cm, right=1.2cm of el] (sp) {\textbf{\small Specificati con}\\[3pt]
    \textbullet\ \textbf{casi d'uso} \textit{(questo capitolo)}\\ \textbullet\ linguaggio naturale\\ \textbullet\ modelli di dominio\\ \textbullet\ mock-up};
  \draw[flow] (el) -- (sp);
  \node[font=\small\bfseries, text=ph1, above=2pt of el] {Requirements Engineering};
\end{tikzpicture}
\caption{Dove si collocano i casi d'uso nella prima fase del ciclo di vita.}
\end{figure}

% ============================================================
\section{Casi d'uso e attori}

\dfn{Casi d'uso (Use Cases)}{
  Modo di descrivere le \textbf{interazioni tra gli utenti e un sistema} usando un \textbf{modello grafico} e un
  \textbf{testo in linguaggio naturale strutturato}. Sono una parte fondamentale di UML (\textit{Unified Modelling Language})
  e servono a rappresentare l'insieme dei \textbf{requisiti funzionali} di un sistema.
}

Un modello dei casi d'uso individua due cose:
\begin{itemize}
  \item gli \textbf{attori}, cioè le \textbf{categorie di utenti} del sistema (non necessariamente esseri umani);
  \item i \textbf{casi d'uso}, cioè i \textbf{tipi di interazione} (le funzionalità) che il sistema offre.
\end{itemize}
Le informazioni aggiuntive su ogni interazione si danno con \textbf{descrizioni testuali} (strutturate) oppure con altri modelli
semi-formali, ad esempio i \textit{Sequence Diagram} o gli \textit{Statechart} di UML.

\info{A cosa servono davvero}{
  I casi d'uso sono prima di tutto uno strumento di \textbf{comunicazione con il cliente} per concordare le funzionalità del
  sistema, quindi devono essere \textbf{il più semplici possibile}. Dicono \textbf{cosa} il sistema deve fare, non \textbf{come}
  sarà implementato: il punto di vista è quello degli utenti, per cui il sistema è una \textbf{scatola nera}.
  Di solito si documentano con un diagramma di alto livello, lo \textbf{Use Case Diagram (UCD)}.
}

\subsection{Gli attori}

\dfn{Attore}{
  Entità \textbf{esterna} che interagisce con il sistema in sviluppo (\textit{System Under Development}, \textbf{SUD}).
  Può essere una classe di utenti umani, un altro sistema, oppure l'ambiente fisico. Si disegna con un \textbf{omino stilizzato}
  e ha un \textbf{nome univoco}.
}

\begin{center}
\begin{tikzpicture}
  \ucactor{a1}{(0,0)}{Customer}
  \ucactor{a2}{(3.2,0)}{Bank Payment Processor}
  \ucactor{a3}{(6.4,0)}{Temperature Sensor}
  \node[note, below=0.1cm of a1-b, align=center] {classe di utenti};
  \node[note, below=0.1cm of a2-b, align=center] {sistema esterno};
  \node[note, below=0.1cm of a3-b, align=center] {ambiente fisico};
\end{tikzpicture}
\end{center}

\warning{{Un attore è un ruolo, non una persona}}{
  Un attore non corrisponde a una singola entità, ma a una \textbf{classe di utenti che hanno lo stesso ruolo}.
  La stessa persona può avere \textbf{ruoli diversi} nello stesso sistema: l'amministratore di un sito web può anche visitarlo
  come utente non autenticato.

  Gli attori sono inoltre \textbf{più grossolani delle personas} del Capitolo 5: a un solo attore possono corrispondere più personas.
}

\Needspace{10\baselineskip}
\textbf{Euristiche per trovare gli attori.} Per individuare gli attori conviene chiedersi:
\begin{enumerate}
  \item quali gruppi di utenti sono \textbf{supportati} dal SUD nel loro lavoro?
  \item quali gruppi di utenti eseguono le \textbf{funzionalità principali} del SUD?
  \item quali gruppi di utenti eseguono le \textbf{funzioni secondarie}, come l'amministrazione?
  \item il SUD interagisce con \textbf{sistemi o software esterni}? Ciascuno di essi sarà un attore.
\end{enumerate}

\subsection{I casi d'uso}

\dfn{Caso d'uso}{
  Funzionalità offerta dal sistema che fornisce un \textbf{beneficio o un'utilità} a uno o più attori.
  Si disegna come un'\textbf{ellisse con un nome} e modella un \textbf{requisito funzionale}.
}

Un caso d'uso non descrive una sola sequenza di azioni: \textbf{astrae molti scenari possibili} per la stessa funzionalità.
Ci possono essere, ad esempio, più modi di iscriversi a un corso.

\begin{center}
\begin{tikzpicture}
  \node[uc, minimum width=3.4cm] (u) at (0,0) {Enroll in a course};
  \foreach \i/\y in {1/0.9, 2/0, 3/-0.9}{
    \node[subchip=ph6, minimum width=2.6cm] (s\i) at (5.2,\y) {Scenario \i};
    \draw[-{Stealth[length=2mm]}, dashed, swgray] (s\i.west) -- (u.east);
  }
  \node[note, right=0.3cm of s2, text width=3.8cm] {ogni scenario è un'\textbf{istanza} del caso d'uso};
  \node[note, below=0.25cm of u, text width=4cm, align=center] {il caso d'uso è la \textbf{classe} degli scenari che usano la stessa funzionalità};
\end{tikzpicture}
\end{center}

\Needspace{6\baselineskip}
In altre parole il rapporto tra caso d'uso e scenario è lo stesso che c'è tra \textbf{classe e oggetto}: lo scenario è
un'istanza concreta, il caso d'uso è la classe di tutti gli scenari con lo stesso obiettivo.

% ============================================================
\section{Le relazioni di un diagramma dei casi d'uso}

Oltre ad attori e casi d'uso, un UCD contiene diversi tipi di \textbf{relazioni} tra di essi.

\subsection{Associazione}

\dfn{Associazione}{
  Un'\textbf{associazione} tra un attore e un caso d'uso indica che l'attore \textbf{può eseguire} quel caso d'uso.
  Si disegna come una \textbf{linea} che collega l'attore all'ellisse.
}

\subsection{Attori secondari e cardinalità}

Un caso d'uso può essere associato a \textbf{più attori}. Il significato è che quegli attori devono \textbf{collaborare}
in qualche modo per eseguire il caso d'uso. In UML 2.5 le associazioni possono avere una \textbf{cardinalità}.

\begin{esempio}{La cassa di un negozio}
\begin{center}
\begin{tikzpicture}
  \ucactor{cu}{(0,1.1)}{Customer}
  \ucactor{ca}{(0,-1.1)}{Cashier}
  \ucactor{bp}{(7,0)}{Bank Payment Processor}
  \node[uc] (pc) at (3.3,0) {Perform Checkout};
  \draw[assoc] ([xshift=9pt]cu) -- (pc);
  \draw[assoc] ([xshift=9pt]ca) -- (pc);
  \draw[assoc] (pc) -- node[pos=0.8, above, font=\scriptsize] {0..1} ([xshift=-9pt]bp);
\end{tikzpicture}
\end{center}
\begin{itemize}
  \item \textbf{Customer} e \textbf{Cashier} devono cooperare per eseguire il checkout.
  \item \textbf{Opzionalmente} (cardinalità \texttt{0..1}) può essere coinvolto anche un \textbf{Bank Payment Processor},
        un sistema esterno: ad esempio quando il cliente paga con carta di credito o di debito.
\end{itemize}
\end{esempio}

\info{Il diagramma non dice \textit{come} collaborano}{
  Gli UCD di UML \textbf{non hanno strumenti} per specificare in che modo i diversi attori partecipano a un caso d'uso.
  Le modalità di interazione e le responsabilità di ciascuno si specificano nelle \textbf{descrizioni testuali}
  che vedremo più avanti in questo capitolo.
}

\subsection{Generalizzazione tra attori}

\dfn{Generalizzazione tra attori}{
  Si usa quando un attore è un \textbf{sottotipo} di un altro. Notazione e significato sono gli stessi dei diagrammi delle classi
  UML: una \textbf{freccia con la punta vuota} che va dall'attore specializzato a quello generale.
  L'attore specializzato può eseguire \textbf{tutti} i casi d'uso che può eseguire il genitore.
}

\begin{center}
\begin{tikzpicture}
  \ucactor{cu}{(0,1.3)}{Customer}
  \ucactor{rc}{(0,-1.2)}{Regular Customer}
  \node[uc] (bi) at (4,1.5) {Buy Items};
  \node[uc] (fp) at (4,-1.0) {Choose fidelity prize};
  \draw[assoc] ([xshift=9pt]cu) -- (bi);
  \draw[assoc] ([xshift=9pt]rc) -- (fp);
  \draw[gen] (0,-0.5) -- (0,0.45);
  \node[note, right=0.5cm of fp, text width=4.4cm] {Regular Customer può scegliere il premio fedeltà \textbf{e} anche comprare
     articoli, perché eredita i casi d'uso di Customer.};
\end{tikzpicture}
\end{center}

\subsection{La relazione «include»}

\dfn{Relazione «include»}{
  Si usa quando due o più casi d'uso hanno una \textbf{parte di comportamento in comune}. La parte comune viene estratta in un
  caso d'uso separato, che viene \textbf{incluso} da tutti i casi d'uso base che la condividono.
  Si disegna come una \textbf{freccia tratteggiata} con lo stereotipo «include», che va \textbf{dal caso d'uso base a quello incluso}.
}

\begin{center}
\begin{tikzpicture}
  \ucactor{cu}{(0,0)}{Customer}
  \node[uc] (wm) at (3.2,0.9) {Withdraw money};
  \node[uc] (dc) at (3.2,-0.9) {Deposit check};
  \node[uc] (ci) at (7.6,0) {Card identification};
  \draw[assoc] ([xshift=9pt]cu) -- (wm);
  \draw[assoc] ([xshift=9pt]cu) -- (dc);
  \draw[dep] (wm) -- node[stereo, above, sloped] {«include»} (ci);
  \draw[dep] (dc) -- node[stereo, below, sloped] {«include»} (ci);
\end{tikzpicture}
\end{center}

Poiché «include» serve soprattutto a \textbf{riusare} parti comuni, quello che resta nel caso d'uso base di solito \textbf{non è
completo da solo}: ha senso solo insieme alle parti incluse. Lo dice anche il verso della freccia: il caso d'uso base
\textbf{dipende} da quello incluso, non viceversa (prelevare senza identificare la carta non ha senso).

\Needspace{6\baselineskip}
La relazione «include» è utile per:
\begin{itemize}
  \item \textbf{decomporre} un'interazione complessa in interazioni più piccole e gestibili;
  \item \textbf{fattorizzare} sequenze di passi comuni a più casi d'uso.
\end{itemize}

\subsection{La relazione «extend»}

\dfn{Relazione «extend»}{
  Si usa quando c'è un \textbf{comportamento aggiuntivo} che \textbf{può} essere aggiunto, eventualmente \textbf{sotto condizione},
  al comportamento di uno o più casi d'uso. Si disegna come una \textbf{freccia tratteggiata} con lo stereotipo «extend», che va
  \textbf{dal caso d'uso che estende a quello esteso}.
}

\begin{center}
\begin{tikzpicture}
  \ucactor{cu}{(0,0)}{Customer}
  \node[uc] (wm) at (3.2,0.9) {Withdraw money};
  \node[uc] (dc) at (3.2,-0.9) {Deposit check};
  \node[uc] (ci) at (7.6,0) {Card identification};
  \node[uc] (pr) at (7.6,1.9) {Print receipt};
  \draw[assoc] ([xshift=9pt]cu) -- (wm);
  \draw[assoc] ([xshift=9pt]cu) -- (dc);
  \draw[dep] (wm) -- node[stereo, below, sloped] {«include»} (ci);
  \draw[dep] (dc) -- node[stereo, below, sloped] {«include»} (ci);
  \draw[dep] (pr) -- node[stereo, above, sloped] {«extend»} (wm);
\end{tikzpicture}
\end{center}

Qui la dipendenza è \textbf{rovesciata} rispetto a «include»:
\begin{itemize}
  \item il caso d'uso \textbf{esteso} (Withdraw money) è definito \textbf{indipendentemente} da quello che lo estende e ha senso
        anche da solo;
  \item il caso d'uso che \textbf{estende} (Print receipt) definisce un comportamento che, di solito, \textbf{non ha senso da solo}
        (stampare la ricevuta di cosa?).
\end{itemize}

\Needspace{12\baselineskip}
\textbf{Extension points.} Gli \textit{extension point} indicano il \textbf{punto preciso} del caso d'uso base in cui può essere
inserito il comportamento di un caso d'uso di estensione. Sono utili quando un caso d'uso può essere esteso in \textbf{più punti}.
Si elencano dentro l'ellisse, sotto il nome, e l'etichetta della freccia «extend» dice quale punto si sta usando.

\begin{center}
\begin{tikzpicture}
  \ucactor{cu}{(0,0)}{Customer}
  \node[uc, minimum width=3.6cm, minimum height=1.35cm] (wm) at (3.4,1.0) {\textbf{Withdraw money}\\[-2pt]
     \rule{2.6cm}{0.4pt}\\[-1pt] \textit{Extension points:}\\ print operation receipt};
  \node[uc] (dc) at (3.4,-1.0) {Deposit check};
  \node[uc] (ci) at (8.2,-0.1) {Card identification};
  \node[uc] (pr) at (9.6,1.6) {Print receipt};
  \draw[assoc] ([xshift=9pt]cu) -- (wm);
  \draw[assoc] ([xshift=9pt]cu) -- (dc);
  \draw[dep] (wm) -- node[stereo, below, sloped] {«include»} (ci);
  \draw[dep] (dc) -- node[stereo, below, sloped] {«include»} (ci);
  \draw[dep] (pr) -- node[stereo, above, align=center] {«extend»\\(print operation receipt)} (wm);
\end{tikzpicture}
\end{center}

\subsection{Generalizzazione tra casi d'uso}

\dfn{Generalizzazione tra casi d'uso}{
  Si usa quando un caso d'uso è una \textbf{specializzazione} di un altro. La notazione è quella standard di UML per la
  specializzazione (freccia con punta vuota verso il caso d'uso generale).
}

\begin{center}
\begin{tikzpicture}
  \node[uc, minimum width=3.2cm] (ci) at (0,0) {\textit{Card identification}};
  \node[uc, minimum width=3.6cm] (g1) at (-4.4,-1.8) {Card identification\\with PIN code};
  \node[uc, minimum width=3.6cm] (g2) at (0,-1.8) {Card identification\\with fingerprint};
  \node[uc, minimum width=3.6cm] (g3) at (4.4,-1.8) {Card identification\\with face recognition};
  \foreach \g in {g1,g2,g3}{ \draw[gen] (\g) -- (ci); }
\end{tikzpicture}
\end{center}

\warning{Generalizzazione o «extend»?}{
  A differenza di «extend», nella generalizzazione \textbf{non ci sono punti precisi} in cui i casi d'uso specializzati si
  discostano dal genitore: le specializzazioni possono essere anche \textbf{molto diverse} dal caso d'uso generale.
  Con «extend» il caso base resta lo stesso e gli si \textbf{aggiunge} qualcosa in un punto preciso; con la generalizzazione il
  caso specializzato \textbf{sostituisce} il genitore.
}

\subsection{Riepilogo delle relazioni}

\begin{table}[H]
\centering
\renewcommand{\arraystretch}{1.5}
\begin{tabularx}{\textwidth}{>{\bfseries}P{2.6cm} >{\centering\arraybackslash}m{4.2cm} L}
\toprule
Relazione & Notazione & Significato \\
\midrule
Associazione &
\tikz[baseline=-0.5ex]{\ucactor{x}{(0,-0.1)}{User}\node[uc, minimum width=0.8cm, minimum height=0.5cm] (b) at (2.4,0) {A};
  \draw[assoc] ([xshift=9pt]x) -- (b);} &
L'attore è associato al caso d'uso: \textit{User può eseguire A}. \\
Include &
\tikz[baseline=-0.5ex]{\node[uc, minimum width=0.8cm, minimum height=0.5cm] (a) at (0,0) {A};
  \node[uc, minimum width=0.8cm, minimum height=0.5cm] (b) at (2.6,0) {B};
  \draw[dep] (a) -- node[stereo, above] {«include»} (b);} &
Il comportamento comune viene fattorizzato: \textit{A include i passi definiti in B}. \\
Extend &
\tikz[baseline=-0.5ex]{\node[uc, minimum width=0.8cm, minimum height=0.5cm] (a) at (0,0) {A};
  \node[uc, minimum width=0.8cm, minimum height=0.5cm] (b) at (2.6,0) {B};
  \draw[dep] (b) -- node[stereo, above] {«extend»} (a);} &
Individua variazioni (eventualmente opzionali): \textit{A può includere la sequenza di passi di B}. \\
Generalizzazione &
\tikz[baseline=-0.5ex]{\node[uc, minimum width=0.8cm, minimum height=0.5cm] (a) at (0,0) {A};
  \node[uc, minimum width=0.8cm, minimum height=0.5cm] (b) at (2.6,0) {B};
  \draw[gen] (b) -- (a);} &
Vale per attori e casi d'uso: \textit{B eredita il comportamento di A}, e \textit{A può essere sostituito da un'istanza di B}. \\
\bottomrule
\end{tabularx}
\caption{Le quattro relazioni di un UCD. Attenzione al verso delle frecce: «include» va dal base all'incluso,
«extend» dall'estensione al caso esteso.}
\end{table}

% ============================================================
\section{Leggere un diagramma dei casi d'uso}

\begin{esempio}{Abbonamenti e buoni regalo}
\begin{center}
\begin{tikzpicture}
  \ucactor{ma}{(-3.8,2.6)}{Manager}
  \ucactor{cu}{(-1.2,2.6)}{Customer}
  \ucactor{ad}{(3.6,2.6)}{Admin}
  \ucactor{ps}{(-1.2,-3.6)}{Payment Service}
  \node[uc] (bs) at (-3.2,0.4) {Buy subscription};
  \node[uc] (bg) at (0.4,0.4) {Buy gift card};
  \node[uc] (vr) at (3.9,0.4) {View sales report};
  \node[uc] (nl) at (-4.4,-1.6) {Sign Up for Newsletter};
  \node[uc] (co) at (-1.2,-1.6) {Checkout};
  \draw[gen] (3.6,3.35) to[bend right=20] (-3.6,3.35);
  \draw[assoc] (ma-b) -- (bs);
  \draw[assoc] (cu-b) -- (bs);
  \draw[assoc] (cu-b) -- (bg);
  \draw[assoc] (ad-b) -- (vr);
  \draw[assoc] (co) -- (-1.2,-2.75);
  \draw[dep] (nl) -- node[stereo, left] {«extend»} (bs);
  \draw[dep] (bs) -- node[stereo, sloped, above] {«include»} (co);
  \draw[dep] (bg) -- node[stereo, sloped, above] {«include»} (co);
\end{tikzpicture}
\end{center}
Il diagramma si legge così:
\begin{itemize}
  \item i \textbf{Customer} possono usare il sistema per comprare un \textbf{buono regalo} o un \textbf{abbonamento};
  \item entrambi i casi d'uso \textbf{includono} un comportamento comune di \textbf{Checkout} (la sequenza di passi del pagamento),
        che richiede la collaborazione di un servizio esterno, il \textbf{Payment Service};
  \item per comprare un abbonamento serve anche la collaborazione di un \textbf{Manager} (o di un Admin): ad esempio perché
        deve approvare l'acquisto;
  \item \textbf{opzionalmente}, mentre compra un abbonamento, il cliente può \textbf{iscriversi alla newsletter} («extend»);
  \item gli \textbf{Admin} sono gli unici che possono vedere i \textbf{report delle vendite}; poiché Admin specializza Manager,
        possono anche fare tutto ciò che fa un Manager.
\end{itemize}
\end{esempio}

% ============================================================
\section{Errori tipici}

\Needspace{16\baselineskip}
\subsection{Errore I: casi d'uso che non sono casi d'uso}

Un caso d'uso deve portare un \textbf{beneficio} all'attore, aiutarlo a completare il suo lavoro o a raggiungere un obiettivo.
Per questo, di solito:
\begin{itemize}
  \item il nome di un caso d'uso contiene un \textbf{verbo};
  \item il nome di un attore è un \textbf{sostantivo}.
\end{itemize}

\begin{center}
\begin{tikzpicture}
  \ucactor{cu}{(0,1.6)}{Customer}
  \node[uc] (w) at (-5.4,0) {Withdraw money};
  \node[uc] (v) at (-2.2,0) {View balance};
  \node[uc=swred] (e) at (1.8,0) {View wrong\\PIN error};
  \node[uc=swred] (p) at (5.4,0) {Press "Deposit"\\button};
  \foreach \u in {w,v,e,p}{ \draw[assoc] (cu-b) -- (\u); }
  \node[font=\scriptsize\bfseries, text=swred, align=center] at (4.6,2.0) {Questi non sono casi d'uso\\per il nostro ATM!};
  \draw[-{Stealth[length=2mm]}, swred, line width=0.8pt] (3.9,1.6) -- (e.north east);
  \draw[-{Stealth[length=2mm]}, swred, line width=0.8pt] (5.2,1.6) -- (p.north);
\end{tikzpicture}
\end{center}

\Needspace{4\baselineskip}
Vedere un messaggio di errore o premere un pulsante non sono obiettivi del cliente: sono \textbf{passi} (o eccezioni) all'interno
di un caso d'uso, e troveranno posto nella sua descrizione testuale.

\subsection{Errore II: più attori sullo stesso caso d'uso}

\error{Due attori sulla stessa ellisse non sono un «oppure»}{
  Se due attori sono associati allo stesso caso d'uso (con cardinalità non nulla), significa che \textbf{entrambi} sono coinvolti,
  e devono \textbf{collaborare}, in \textbf{ogni istanza} (scenario) di quel caso d'uso. Nell'esempio della cassa, Customer e Cashier
  partecipano insieme a ogni checkout.

  \textbf{Non} significa che ciascuno dei due può eseguire il caso d'uso per conto proprio! Per dire «lo possono fare entrambi»
  si usa invece una generalizzazione tra attori.
}

\subsection{Errore III: generalizzazioni inutili}

\begin{center}
\begin{tikzpicture}
  \ucactor{em}{(0,1.8)}{Employee}
  \ucactor{mg}{(0,-0.4)}{Manager}
  \ucactor{te}{(0,-2.7)}{Technical Employee}
  \node[uc] (ot) at (4,2.0) {Open Ticket};
  \node[uc] (mt) at (4,-2.0) {Manage Ticket};
  \node[uc] (ct) at (4,-3.1) {Close Ticket};
  \draw[assoc] ([xshift=9pt]em) -- (ot);
  \draw[assoc] ([xshift=9pt]te) -- (mt);
  \draw[assoc] ([xshift=9pt]te) -- (ct);
  \draw[gen] (0,0.3) -- (0,0.95);
  \draw[gen] (-0.35,-2.6) to[bend left=45] (-0.35,2.0);
  \node[font=\scriptsize\bfseries, text=swred, anchor=west] at (0.9,-0.25) {Manager non ha casi d'uso propri};
\end{tikzpicture}
\end{center}

\warning{Ogni attore specializzato deve avere i suoi casi d'uso}{
  Un attore specializzato può già eseguire tutti i casi d'uso dei suoi antenati. Se \textbf{non ha casi d'uso propri}
  (come Manager qui sopra), allora o la generalizzazione è \textbf{inutile}, oppure \textbf{mancano} dei casi d'uso nel diagramma.
}

\subsection{Errore IV: «include» per le relazioni temporali}

\begin{center}
\begin{tikzpicture}
  \ucactor{cu}{(0,0)}{Customer}
  \node[uc] (vc) at (3.2,0.8) {View Catalogue};
  \node[uc] (bi) at (3.2,-0.8) {Buy Items};
  \node[uc=swred] (pl) at (7.6,0) {Perform Login};
  \draw[assoc] ([xshift=9pt]cu) -- (vc);
  \draw[assoc] ([xshift=9pt]cu) -- (bi);
  \draw[dep] (vc) -- node[stereo, above, sloped] {«include»} (pl);
  \draw[dep] (bi) -- node[stereo, below, sloped] {«include»} (pl);
\end{tikzpicture}
\end{center}

\error{«include» non vuol dire «prima»}{
  Questo diagramma dice che il Customer deve eseguire \textbf{tutti i passi} di Perform Login \textbf{ogni volta} che vuole
  comprare articoli o consultare il catalogo! Il fatto che l'utente debba essere autenticato è una \textbf{relazione temporale},
  e va espresso con una \textbf{precondizione} nella descrizione testuale del caso d'uso («il cliente ha effettuato il login»).
}

\subsection{Errore IV bis: diagrammi troppo complessi}

Un diagramma dei casi d'uso non deve diventare complesso e disordinato:
\begin{itemize}
  \item le relazioni tra casi d'uso e le generalizzazioni tra attori vanno usate \textbf{con moderazione};
  \item la modellazione deve restare a un livello di \textbf{astrazione sufficientemente alto};
  \item bisogna essere \textbf{ordinati}: ad esempio evitare le linee che si intersecano.
\end{itemize}

\error{Un diagramma complesso è il segno di un cattivo analista}{
  Un UCD pieno di ellissi, frecce «include»/«extend» incrociate e gerarchie di attori non comunica niente al cliente,
  che è proprio il suo destinatario. I dettagli vanno nelle descrizioni testuali, non nel diagramma.
}

% ============================================================
\section{Specificare i casi d'uso}

Lo UCD dà una panoramica \textbf{di altissimo livello} dei requisiti funzionali: non è abbastanza dettagliato per stabilire i
\textbf{requisiti di sistema}. Per \textbf{ogni} caso d'uso del diagramma serve quindi una \textbf{specifica dettagliata}, che descriva
ogni aspetto dell'interazione dal punto di vista dell'attore, compresi tutti gli scenari e le variazioni possibili.

\Needspace{10\baselineskip}
Come per gli scenari del Capitolo 5, la descrizione testuale di un caso d'uso comprende in genere:
\begin{enumerate}
  \item cosa il sistema e gli utenti si \textbf{aspettano} quando il caso d'uso inizia;
  \item il \textbf{flusso normale} degli eventi (\textbf{main scenario});
  \item cosa può causare \textbf{errori} e come si gestiscono i problemi che ne derivano;
  \item lo \textbf{stato del sistema} al termine del caso d'uso.
\end{enumerate}

\subsection{I formati}

I casi d'uso si possono scrivere con diversi livelli di formalità.

\begin{table}[H]
\centering
\renewcommand{\arraystretch}{1.4}
\begin{tabularx}{\textwidth}{>{\bfseries\leavevmode\color{ph4}}P{3.2cm} L P{4.4cm}}
\toprule
Formato & Descrizione & Quando si usa \\
\midrule
Brief & Riassunto conciso di un paragrafo, di solito del solo scenario principale di successo. &
Nelle \textbf{prime fasi} della specifica, per farsi un'idea veloce di argomento e perimetro. \\
Casual & Paragrafi informali: più paragrafi che coprono i vari scenari. & Come il formato brief: nelle \textbf{prime fasi}. \\
Fully-dressed & Tutti i passi e tutte le variazioni sono scritti in dettaglio, con sezioni di supporto come precondizioni
e garanzie di successo. & \textbf{Più avanti}, come \textbf{base per il contratto} e per specificare in dettaglio il comportamento
del sistema. \\
\bottomrule
\end{tabularx}
\end{table}

\subsection{Il template di Cockburn}

Per le descrizioni fully-dressed sono stati proposti molti formati; noi usiamo un template basato su quello di
\textbf{Alistair Cockburn}. I nomi dei campi restano in inglese.

\begin{table}[H]
\centering
\renewcommand{\arraystretch}{1.35}
\small
\begin{tabularx}{\textwidth}{>{\bfseries}P{3.3cm} P{1.3cm} L L L}
\toprule
\rowcolor{ph4!12} USE CASE \#X & \multicolumn{4}{l}{\textbf{Nome del caso d'uso}} \\
Goal in Context & \multicolumn{4}{l}{Obiettivo del caso d'uso} \\
Preconditions & \multicolumn{4}{l}{Tutte le condizioni che devono valere per poter iniziare il caso d'uso} \\
Success End Condition & \multicolumn{4}{l}{Stato del sistema se il caso d'uso ha successo} \\
Failed End Condition & \multicolumn{4}{l}{Stato del sistema se il caso d'uso fallisce} \\
Primary Actor & \multicolumn{4}{l}{Attore principale del caso d'uso} \\
Trigger & \multicolumn{4}{l}{Azione dell'attore principale che avvia il caso d'uso} \\
\midrule
Main Scenario & \textbf{Step} & \textbf{Actor 1} & \textbf{Actor n} & \textbf{System} \\
 & 1 & Azione di avvio (trigger) & & \\
 & 2 & & & Risposta \\
 & \ldots & Azione 2 & \ldots & \ldots \\
 & n & & & Azione finale \\
\midrule
Extension \#1\newline {\normalfont\footnotesize(breve descrizione)} & \textbf{Step} & \textbf{Actor 1} & \textbf{Actor n} & \textbf{System} \\
 & x \textless condizione\textgreater & \ldots & \ldots & \ldots \\
 & \ldots & & & Azione finale (eventualmente torna a un passo del main scenario) \\
\midrule
Extension \#n & \multicolumn{4}{l}{\ldots} \\
\midrule
Open Issues & \multicolumn{4}{l}{Aspetti ancora da chiarire. \textbf{Alla consegna del documento deve essere vuoto.}} \\
\bottomrule
\end{tabularx}
\caption{Il template fully-dressed di Cockburn.}
\end{table}

\subsection{Main scenario ed estensioni}

\dfn{Main scenario}{
  Sequenza di azioni che avviene quando nel caso d'uso \textbf{tutto va come previsto}.
}

Ma un caso d'uso si può eseguire in \textbf{modi diversi} (l'utente può autenticarsi con il PIN oppure con l'impronta digitale)
e in qualunque punto può verificarsi un \textbf{errore}. Quando si definisce il comportamento funzionale del sistema bisogna
descrivere anche queste sequenze alternative.

\dfn{Extension}{
  Sequenza di azioni \textbf{alternativa} al main scenario, che parte da un passo preciso quando si verifica una
  \textbf{condizione}. Termina il caso d'uso oppure \textbf{torna} a un passo del main scenario.
}

\info{Dove sta davvero il lavoro}{
  Di solito c'è \textbf{molto più testo nelle estensioni} che nel main scenario: il percorso «felice» è uno solo,
  i modi in cui qualcosa può andare storto o diversamente sono tanti.
}

% ============================================================
\section{Un esempio completo: il bancomat}

Vogliamo descrivere in formato fully-dressed il caso d'uso \textbf{Withdraw money} dello UCD del bancomat (ATM) visto sopra.

\Needspace{24\baselineskip}
\subsection{I mock-up}

Prima di scrivere i passi è utile disegnare dei \textbf{mock-up} delle schermate (Capitolo 5): i passi del caso d'uso potranno
farvi riferimento per nome (M1, M2, \ldots).

\newcommand{\atmscreen}[3]{%
  \begin{scope}[shift={#1}]
    \draw[rounded corners=3pt, line width=0.9pt, fill=white, draw=black!70] (0,0) rectangle (3.6,2.3);
    \node[chip=ph4, anchor=north west] at (0.08,2.22) {#2};
    #3
  \end{scope}}
\newcommand{\atmbtn}[3]{\node[draw=black!60, fill=swlight, rounded corners=1pt, font=\tiny, minimum width=#3, minimum height=0.3cm, inner sep=1pt] at #1 {#2};}

\begin{figure}[H]
\centering
\begin{tikzpicture}
  \atmscreen{(0,0)}{M1}{\node[font=\scriptsize, align=center] at (1.8,1.05) {Insert your ATM Card\\to begin};}
  \atmscreen{(4.4,0)}{M2}{
     \node[font=\scriptsize\bfseries, align=center] at (1.95,1.95) {Select an Authentication\\Method};
     \atmbtn{(0.65,0.75)}{PIN}{0.8cm}
     \atmbtn{(1.8,0.75)}{Fingerprint}{0.8cm}
     \atmbtn{(2.95,0.75)}{Face}{0.8cm}}
  \atmscreen{(8.8,0)}{M3}{
     \node[font=\scriptsize\bfseries] at (2.0,2.0) {Insert your PIN};
     \node[draw=black!60, font=\scriptsize, minimum width=1.1cm] at (0.8,1.05) {****};
     \foreach \r/\row in {0/{7,8,9},1/{4,5,6},2/{1,2,3}}{
       \foreach \c [count=\k from 0] in \row { \atmbtn{(1.9+\k*0.38,1.55-\r*0.36)}{\c}{0.32cm} } }
     \atmbtn{(2.28,0.47)}{0}{0.32cm}
     \atmbtn{(3.25,1.55)}{DEL}{0.45cm}
     \atmbtn{(3.25,0.91)}{OK}{0.45cm}}
  \draw[flow] (3.65,1.15) -- (4.35,1.15);
  \draw[flow] (8.05,1.15) -- (8.75,1.15);
  \draw[flow] (10.6,-0.05) -- (10.6,-0.6);
  \atmscreen{(8.8,-2.95)}{M4}{
     \node[font=\scriptsize\bfseries] at (2.05,2.0) {Select an Operation};
     \atmbtn{(1.8,1.2)}{WITHDRAW}{2.2cm}
     \atmbtn{(1.8,0.6)}{DEPOSIT CHECK}{2.2cm}}
  \atmscreen{(4.4,-2.95)}{M5}{
     \node[font=\scriptsize\bfseries] at (2.05,2.0) {Select Amount};
     \foreach \v [count=\k from 0] in {20,50,100}{ \atmbtn{(0.7+\k*1.1,1.4)}{\v\ €}{0.85cm} }
     \foreach \v [count=\k from 0] in {150,250,500}{ \atmbtn{(0.7+\k*1.1,0.9)}{\v\ €}{0.85cm} }
     \draw (0.3,0.3) rectangle (0.48,0.48); \node[font=\tiny, anchor=west] at (0.52,0.39) {Receipt};}
  \atmscreen{(0,-2.95)}{M6}{\node[font=\scriptsize, align=center] at (1.8,1.05) {Collect your money\\and your card};}
  \draw[flow] (8.75,-1.8) -- (8.05,-1.8);
  \draw[flow] (4.35,-1.8) -- (3.65,-1.8);
\end{tikzpicture}
\caption{I mock-up M1--M6 del prelievo al bancomat, nell'ordine del main scenario.}
\end{figure}

\subsection{Il caso d'uso fully-dressed}

\begin{table}[H]
\centering
\renewcommand{\arraystretch}{1.3}
\small
\begin{tabularx}{\textwidth}{>{\bfseries}P{3.3cm} P{1.2cm} L L}
\toprule
\rowcolor{ph4!12} USE CASE \#1 & \multicolumn{3}{l}{\textbf{Withdraw money}} \\
Goal in Context & \multicolumn{3}{l}{Un cliente vuole prelevare denaro dal bancomat.} \\
Preconditions & \multicolumn{3}{p{11.2cm}}{Il cliente ha un conto presso la banca e possiede una carta bancomat. Ha 500 € sul conto.} \\
Success End Condition & \multicolumn{3}{p{11.2cm}}{Il sistema registra l'operazione di prelievo ed eroga il denaro richiesto.} \\
Failed End Condition & \multicolumn{3}{l}{Non viene effettuata alcuna transazione.} \\
Primary Actor & \multicolumn{3}{l}{Customer} \\
Trigger & \multicolumn{3}{p{11.2cm}}{Il cliente si avvicina al bancomat e tocca lo schermo per attivarlo.} \\
\midrule
Main Scenario & \textbf{Step} & \textbf{Customer} & \textbf{System} \\
 & 1 & Tocca lo schermo & \\
 & 2 & & Mostra M1 \\
 & 3 & Inserisce la carta & \\
 & 4 & & Mostra M2 \\
 & 5 & Tocca il pulsante «PIN» & \\
 & 6 & & Mostra M3 \\
 & 7 & Inserisce il PIN & \\
 & 8 & & Mostra M4 \\
 & 9 & Tocca il pulsante «Withdraw» & \\
 & 10 & & Mostra M5 \\
 & 11 & Tocca il pulsante «50 €» & \\
 & 12 & & Eroga il denaro, espelle la carta, mostra M6 \\
\bottomrule
\end{tabularx}
\end{table}

Si noti la struttura a \textbf{ping-pong}: ogni azione del cliente è seguita da una risposta del sistema, e ogni risposta del sistema
fa riferimento a un mock-up preciso.

\Needspace{12\baselineskip}
\subsection{Le estensioni}

Ora ci si chiede sistematicamente \textbf{cosa può andare storto} e \textbf{cosa può andare diversamente}.

\begin{center}
\begin{tikzpicture}
  \node[card=swred, text width=7.2cm, align=left, minimum height=3.2cm] (w) {\textbf{Cosa può andare storto?}\\[3pt]\small
    \textbullet\ il PIN non è corretto\\ \textbullet\ il cliente non ha abbastanza soldi sul conto\\
    \textbullet\ il bancomat non ha abbastanza contanti\\ \textbullet\ la carta è segnalata come rubata\\
    \textbullet\ la carta è illeggibile\\ \textbullet\ \ldots};
  \node[card=ph6, text width=7.2cm, align=left, minimum height=3.2cm, right=0.4cm of w] {\textbf{Cosa può andare diversamente?}\\[3pt]\small
    \textbullet\ il cliente si autentica con l'impronta digitale o con il riconoscimento facciale\\
    \textbullet\ il cliente chiede la ricevuta stampata};
\end{tikzpicture}
\end{center}

\Needspace{6\baselineskip}
Ognuno di questi scenari va descritto con un'estensione. Il passo di un'estensione si numera con il \textbf{numero del passo del
main scenario} da cui devia, seguito da una \textbf{lettera} che identifica l'estensione (7a, 8a, \ldots\ per la prima; 11b, 12b,
\ldots\ per la seconda). Il primo passo riporta tra parentesi angolari la \textbf{condizione} che attiva l'estensione.

\begin{figure}[H]
\centering
\begin{minipage}[c]{0.66\textwidth}
\renewcommand{\arraystretch}{1.3}\small
\begin{tabularx}{\linewidth}{>{\bfseries}P{2.4cm} P{1.9cm} L L}
\toprule
Extension \#1\newline{\normalfont\footnotesize(il cliente inserisce un PIN non valido)} & \textbf{Step} & \textbf{Customer} & \textbf{System} \\
 & 7a \textless PIN errato\textgreater & Inserisce il PIN & \\
 & 8a & & Mostra M7 e termina il caso d'uso \\
\bottomrule
\end{tabularx}
\end{minipage}\hfill
\begin{minipage}[c]{0.3\textwidth}
\centering
\begin{tikzpicture}
  \atmscreen{(0,0)}{M7}{\node[font=\scriptsize, align=center] at (1.8,1.0) {Invalid PIN -- Authentication\\denied.\\Recover your card\\from the tray.};}
\end{tikzpicture}
\end{minipage}
\end{figure}

\begin{figure}[H]
\centering
\begin{minipage}[c]{0.66\textwidth}
\renewcommand{\arraystretch}{1.3}\small
\begin{tabularx}{\linewidth}{>{\bfseries}P{2.4cm} P{1.9cm} L L}
\toprule
Extension \#2\newline{\normalfont\footnotesize(il cliente non ha abbastanza soldi)} & \textbf{Step} & \textbf{Customer} & \textbf{System} \\
 & 11b & Tocca il pulsante «500 €» & \\
 & 12b & & Mostra M8 \\
 & 13b & Clicca «OK» & \\
 & 14b & & Torna al passo 10 del main scenario \\
\bottomrule
\end{tabularx}
\end{minipage}\hfill
\begin{minipage}[c]{0.3\textwidth}
\centering
\begin{tikzpicture}
  \atmscreen{(0,0)}{M8}{\node[font=\scriptsize, align=center] at (1.8,1.25) {Your balance is smaller than\\the amount you are trying\\to withdraw.\\Select a smaller amount.};
     \atmbtn{(1.8,0.3)}{OK}{0.7cm}}
\end{tikzpicture}
\end{minipage}
\end{figure}

\info{{Estensioni che terminano, estensioni che rientrano}}{
  Le due estensioni mostrano i due modi in cui può finire una deviazione: l'Extension \#1 \textbf{termina} il caso d'uso
  (si arriva alla \textit{Failed End Condition}: nessuna transazione), l'Extension \#2 invece \textbf{rientra} nel main scenario
  al passo 10, dove il cliente può scegliere un importo più piccolo.
}

\warning{Precondizioni ed estensioni devono essere coerenti}{
  Nella slide la precondizione dice che il cliente ha 500 € sul conto, ma l'Extension \#2 scatta proprio quando il cliente prova a
  prelevare 500 € e il saldo è insufficiente. Le due cose non possono valere insieme: in un caso d'uso reale la precondizione
  andrebbe riformulata (ad esempio «il cliente ha almeno 50 € sul conto», quanto serve al main scenario). È un buon esempio del
  tipo di incoerenza da cercare in fase di \textbf{validazione}.
}

\info{Per approfondire}{
  Oltre al libro originale di Cockburn esistono molte linee guida per scrivere casi d'uso efficaci, ad esempio
  B.~Wei, L.~Deng, Y.~Wang, \textit{A Practical Style Guide and Templates Repository for Writing Effective Use Cases}, SERA 2022
  (\href{https://doi.org/10.1007/978-3-031-09145-2_5}{doi:10.1007/978-3-031-09145-2\_5}). L'articolo propone anche dei template
  di casi d'uso generici (\url{https://github.com/Washingtonwei/use-case-style-guide}) e un esempio completo,
  il \textit{SuperFrog Scheduler}.
}

% ============================================================
\section{Esercizi}

\begin{esercizio}{Include, extend o generalizzazione?}
Per ciascuna coppia di casi d'uso indica la relazione più adatta e il verso della freccia.
\begin{enumerate}
  \item \textit{Place order} e \textit{Apply discount code} (il codice sconto si inserisce solo se il cliente ne ha uno).
  \item \textit{Place order}, \textit{Book table} e \textit{Pay} (entrambi richiedono sempre un pagamento).
  \item \textit{Pay} e \textit{Pay with PayPal}.
\end{enumerate}

\textbf{Soluzione.}
(1) «extend», da \textit{Apply discount code} a \textit{Place order}: è comportamento opzionale e l'ordine ha senso anche senza.
(2) «include», da \textit{Place order} e da \textit{Book table} verso \textit{Pay}: è una parte comune, sempre eseguita, fattorizzata.
(3) Generalizzazione, da \textit{Pay with PayPal} a \textit{Pay}: è una specializzazione che può sostituire il genitore.
\end{esercizio}

\begin{esercizio}{Trova gli errori}
Un analista ha disegnato per un sistema di e-commerce uno UCD con l'attore \textit{Customer} associato ai casi d'uso
\textit{Search product}, \textit{Click "Add to cart" button} e \textit{Buy items}; \textit{Buy items} «include» \textit{Perform login}.
Inoltre \textit{Customer} e \textit{Admin} sono entrambi associati a \textit{Update product catalogue}, perché «possono farlo tutti e due».
Individua gli errori.

\textbf{Soluzione.}
\textit{Click "Add to cart" button} è un passo dell'interfaccia, non un caso d'uso (Errore I).
«include» verso \textit{Perform login} esprime una relazione temporale: va sostituito da una precondizione (Errore IV).
Associare due attori allo stesso caso d'uso significa che devono \textbf{collaborare} in ogni scenario, non che possono farlo
entrambi (Errore II); qui andrebbe associato solo l'attore autorizzato.
\end{esercizio}

\Needspace{26\baselineskip}
\gbox{In sintesi}{azure-gradient-6!80!black}{
\begin{itemize}
  \item I \textbf{casi d'uso} (UML) descrivono i requisiti funzionali con un diagramma (UCD) e descrizioni testuali: dicono
        \textbf{cosa} fa il sistema, visto come scatola nera.
  \item Gli \textbf{attori} sono ruoli esterni al SUD (utenti, sistemi, ambiente fisico); i \textbf{casi d'uso} sono funzionalità
        che portano un beneficio agli attori e astraggono molti scenari.
  \item Relazioni: \textbf{associazione} (con eventuale cardinalità), \textbf{generalizzazione} di attori e casi d'uso,
        \textbf{«include»} (parte comune, sempre eseguita, dal base all'incluso) e \textbf{«extend»} (comportamento opzionale o
        condizionale, dall'estensione al caso esteso, eventualmente in un \textit{extension point}).
  \item Errori tipici: casi d'uso senza verbo o che sono semplici passi, più attori sullo stesso caso d'uso letti come «oppure»,
        attori specializzati senza casi d'uso propri, «include» per il login, diagrammi troppo complessi.
  \item Ogni caso d'uso va specificato (formati \textbf{brief}, \textbf{casual}, \textbf{fully-dressed}); il template di
        \textbf{Cockburn} prevede goal, precondizioni, condizioni di fine, attore principale, trigger, \textbf{main scenario},
        \textbf{estensioni} e open issues (vuoti alla consegna).
\end{itemize}
}
